\chapter{1960年全国卷}

\section{解答题}

\begin{problem}
\begin{enumerate}
\item
解方程 \(\sqrt{2x^{2}-5}-\sqrt{5x-7}=0\). （限定在实数范围内）

\item
有 \(5\) 组篮球队，每组 \(6\) 队. 首先每组中各队进行单循环比赛（即每两队比赛一次），然后各组的冠军再进行单循环比赛. 问先后共比赛多少次？

\item
求证等比数列的各项的对数组成等差数列. （这里所说的等比数列的各项都是正数）

\item
求使等式 \(\sqrt{1-\sin^{2}\frac{x}{2}}=\cos\frac{x}{2}\) 成立的 \(x\) 值的范围. （\(x\) 是 \(0^{\circ}\) 到 \(720^{\circ}\) 的角）

\item
如图，用钢珠测量机件上一小孔的直径. 所用钢珠的中心是 \(O\)，直径是 \(12\text{ 毫米}\). 将钢珠放在这小孔上，测得钢珠上端 \(D\) 到机件平面的距离 \(CD\) 是 \(9\text{ 毫米}\). 求这小孔的直径 \(AB\) 的长.

\item
四棱锥 \(P-ABCD\) 的底面是一正方形. \(PA\) 与底面垂直. 已知 \(PA\) 是 \(3\text{ 厘米}\)，\(P\) 点到 \(BC\) 的距离是 \(5\text{ 厘米}\). 求 \(PC\) 的长.
\end{enumerate}

\begingroup
\def\theprobnum{5}%
\bitmapfigure[max width=0.34\linewidth]{img/1960/national/q5_fig1.png}
\endgroup
\end{problem}

\begin{solution}
\begin{enumerate}
\item
先由两个平方根都有意义，得
\(2x^{2}-5\geqslant0\)，\(5x-7\geqslant0\).
由第一式得 \(x\leqslant-\sqrt{\frac52}\) 或 \(x\geqslant\sqrt{\frac52}\)，由第二式得 \(x\geqslant\frac75\). 两者合并为 \(x\geqslant\sqrt{\frac52}\). 在这个定义域内，两边均为非负数，原方程等价于
\[
\sqrt{2x^{2}-5}=\sqrt{5x-7}
\Longleftrightarrow 2x^{2}-5=5x-7
\]
整理得 \(2x^{2}-5x+2=0\)，即 \((2x-1)(x-2)=0\)，所以候选值为 \(x=\frac12\) 或 \(x=2\). 其中 \(\frac12\) 不在定义域内，舍去；代入 \(x=2\) 时两被开方式均为 \(3\)，确实成立. 故 \(x=2\).

\item
每组有 \(6\) 队，任取其中两队比赛一次，所以每组比赛 \(\mathrm C_6^2=\frac{6\cdot5}{2}=15\) 次，五组共比赛 \(5\mathrm C_6^2=5\times15=75\) 次. 五个小组冠军再进行单循环比赛，比赛次数为 \(\mathrm C_5^2=\frac{5\cdot4}{2}=10\) 次.
因此先后共比赛 \(75+10=85\) 次.

\item
设等比数列的首项为 \(a>0\)，公比为 \(q>0\)，则第 \(k\) 项为 \(a_k=aq^{k-1}\). 由于各项均为正数，对每个 \(k\) 都有
\[
\log a_k=\log\left(aq^{k-1}\right)=\log a+(k-1)\log q
\]
于是相邻两项对数之差为
\[
\log a_{k+1}-\log a_k=\bigl(\log a+k\log q\bigr)-\bigl(\log a+(k-1)\log q\bigr)=\log q
\]
与 \(k\) 无关. 因此各项的对数构成首项为 \(\log a\)、公差为 \(\log q\) 的等差数列，命题得证.

\item
因为 \(0^{\circ}\leqslant x\leqslant720^{\circ}\)，所以 \(0^{\circ}\leqslant\frac{x}{2}\leqslant360^{\circ}\). 原等式左边为
\[
\sqrt{1-\sin^{2}\frac{x}{2}}=\sqrt{\cos^{2}\frac{x}{2}}=\left|\cos\frac{x}{2}\right|
\]
故原等式成立当且仅当 \(\left|\cos\frac{x}{2}\right|=\cos\frac{x}{2}\)，也就是 \(\cos\frac{x}{2}\geqslant0\). 在 \([0^{\circ},360^{\circ}]\) 内，这等价于
\[
0^{\circ}\leqslant\frac{x}{2}\leqslant90^{\circ}\quad\text{或}\quad270^{\circ}\leqslant\frac{x}{2}\leqslant360^{\circ}
\]
所以 \(0^{\circ}\leqslant x\leqslant180^{\circ}\) 或 \(540^{\circ}\leqslant x\leqslant720^{\circ}\).

\item
解法一：设小孔直径 \(AB=x\text{ 毫米}\). 直径 \(DE\) 垂直于弦 \(AB\)，所以它经过弦的中点 \(C\)，从而 \(AC=\frac{x}{2}\text{ 毫米}\). 又因为 \(DE=12\text{ 毫米}\)、\(CD=9\text{ 毫米}\)，所以
\[
CE=DE-CD=12-9=3\text{ 毫米}
\]
两条弦 \(AB\)、\(DE\) 在圆内相交于 \(C\)，由相交弦定理，
\[
CA\cdot CB=CD\cdot CE
\]
由于 \(CA=CB=\frac{x}{2}\)，得
\(\left(\frac{x}{2}\right)^2=9\times3\)，\(x^2=108\).
长度取正值，所以 \(x=6\sqrt3\text{ 毫米}\). 答：小孔的直径是 \(6\sqrt3\text{ 毫米}\).

解法二：由 \(OD=6\text{ 毫米}\)、\(CD=9\text{ 毫米}\)，且图中 \(C\) 在 \(O\) 与 \(E\) 的同一直线上，得
\[
OC=CD-OD=9-6=3\text{ 毫米}
\]
连接 \(OA\). 因为 \(OC\perp AB\)，且 \(C\) 是弦 \(AB\) 的中点，所以 \(AC=\frac{x}{2}\text{ 毫米}\). 在直角三角形 \(OAC\) 中，
\(OA^2=OC^2+AC^2\)，\(6^2=3^2+\left(\frac{x}{2}\right)^2\).
因此 \(\frac{x^2}{4}=27\)，长度取正值得 \(x=6\sqrt3\text{ 毫米}\). 答：小孔的直径是 \(6\sqrt3\text{ 毫米}\).

\item
由 \(PA\perp\) 平面 \(ABCD\)，知 \(PA\perp BC\). 又底面为正方形，故 \(AB\perp BC\). 直线 \(PA\)、\(AB\) 相交并且都在平面 \(PAB\) 内，因此由线面垂直判定的逆定理，\(BC\perp\) 平面 \(PAB\)，从而 \(BC\perp PB\). 所以点 \(P\) 到直线 \(BC\) 的距离为 \(PB=5\text{ 厘米}\).

因为 \(PA\perp AB\)，所以在直角三角形 \(PAB\) 中，
\[
AB^2=PB^2-PA^2=5^2-3^2=16
\]
故 \(AB=4\text{ 厘米}\). 正方形 \(ABCD\) 的边长相等，所以 \(BC=4\text{ 厘米}\). 在直角三角形 \(PBC\) 中，
\[
PC^2=PB^2+BC^2=5^2+4^2=41
\]
因此 \(PC=\sqrt{41}\text{ 厘米}\).
\end{enumerate}
\FigureLayoutDeclare{img/1960/national/q6_solution_fig1.png}{6.9cm}{0bp 0bp 0bp 0bp}
\solutionfigure{\bitmapfigure[max width=0.34\linewidth]{img/1960/national/q6_solution_fig1.png}}
\end{solution}

\begin{problem}
有一直圆柱，高是 \(20\text{ 厘米}\)，底面半径是 \(5\text{ 厘米}\). 它的一个内接长方体的体积是 \(800\text{ 立方厘米}\). 求这长方体的底面的长和宽.
\end{problem}

\begin{solution}
设内接长方体底面的长和宽分别为 \(x\text{ 厘米}\) 和 \(y\text{ 厘米}\)，其中 \(x>0\)、\(y>0\). 长方体底面是圆柱底面内接的矩形，矩形的四个顶点在圆周上，所以它的两条对角线都是圆柱底面的直径；因此对角线长为 \(2\times5=10\text{ 厘米}\)，从而
\[
\begin{cases}
x^{2}+y^{2}=10^{2},\\
20xy=800
\end{cases}
\]
其中第二式来自长方体的高等于圆柱的高 \(20\text{ 厘米}\). 方程组化为
\[
\begin{cases}
x^{2}+y^{2}=100,\\
2xy=80
\end{cases}
\]
两式相加、相减分别得
\((x+y)^2=x^2+y^2+2xy=100+80=180\)，\((x-y)^2=x^2+y^2-2xy=100-80=20\).
因为 \(x+y>0\)，所以 \(x+y=6\sqrt5\)；又因为 \((x-y)^2=20\)，所以 \(x-y=\pm2\sqrt5\).

分别解下面的两个方程组：
\[
\begin{cases}
x+y=6\sqrt5,\\
x-y=2\sqrt5
\end{cases}
\qquad
\begin{cases}
x+y=6\sqrt5,\\
x-y=-2\sqrt5
\end{cases}
\]
分别解这两个方程组，得到
\[
\begin{cases}
x_1=4\sqrt5,\\
y_1=2\sqrt5
\end{cases}
\qquad
\begin{cases}
x_2=2\sqrt5,\\
y_2=4\sqrt5
\end{cases}
\]
两者只交换长、宽位置，均满足原条件. 答：底面的长和宽分别是 \(4\sqrt5\text{ 厘米}\) 和 \(2\sqrt5\text{ 厘米}\).
\end{solution}

\begin{problem}
从一船上看到在它的南 \(30^{\circ}\) 东的海面上有一灯塔. 船以 \(30\text{ 浬/小时}\) 的速度向正东南方向航行半小时后，看到这个灯塔在船的正西方. 问这时船与灯塔的距离是多少？（精确到 \(0.1\text{ 浬}\)）
\end{problem}

\begin{solution}
解法一：如图，\(A\) 是原来船的位置，\(B\) 是灯塔，\(C\) 是航行半小时后船的位置. 船行驶的路程为
\[
AC=30\times\frac12=15\text{ 浬}
\]
从 \(A\) 看，\(AB\) 为南偏东 \(30^{\circ}\) 的方向，\(AC\) 为正东南方向，即南偏东 \(45^{\circ}\)，所以 \(\angle BAC=45^{\circ}-30^{\circ}=15^{\circ}\). 航行后灯塔在船的正西方，因此 \(BC\) 为东西方向；在点 \(B\) 处，\(BA\) 指向北偏西 \(30^{\circ}\)，故 \(\angle ABC=120^{\circ}\). 由正弦定理，
\[
\begin{gathered}
BC=\frac{AC\sin15^{\circ}}{\sin120^{\circ}}=\frac{15\sin15^{\circ}}{\sin60^{\circ}},\\
\sin15^{\circ}=\sin(45^{\circ}-30^{\circ})=\sin45^{\circ}\cos30^{\circ}-\cos45^{\circ}\sin30^{\circ},\\
BC=\frac{15(\sin45^{\circ}\cos30^{\circ}-\cos45^{\circ}\sin30^{\circ})}{\sin60^{\circ}}=\frac{5\sqrt2(3-\sqrt3)}{2}\approx4.482.
\end{gathered}
\]
按题意精确到 \(0.1\text{ 浬}\)，得 \(BC\approx4.5\text{ 浬}\). 答：这时船与灯塔的距离是 \(4.5\text{ 浬}\).
\solutionfigure{\bitmapfigure[max width=0.34\linewidth]{img/1960/national/q3_solution_fig1.png}}

解法二：过 \(A\) 作 \(AD\perp BC\)，与 \(CB\) 的延长线交于 \(D\). 因为 \(AC\) 与正东南方向成 \(45^{\circ}\)，三角形 \(ADC\) 为等腰直角三角形，所以
\(AC=15\text{ 浬}\)，\(DA=DC=AC\sin45^{\circ}=\frac{15\sqrt2}{2}\text{ 浬}\).
又 \(AB\) 为南偏东 \(30^{\circ}\)，所以在直角三角形 \(ADB\) 中，
\[
DB=DA\tan30^{\circ}=\frac{15\sqrt2}{2}\times\frac{\sqrt3}{3}=\frac{5\sqrt6}{2}\text{ 浬}
\]
由图中点的顺序，\(BC=DC-DB\)，故
\[
BC=\frac{15\sqrt2}{2}-\frac{5\sqrt6}{2}=\frac{5(3\sqrt2-\sqrt6)}{2}\approx4.482\text{ 浬}\approx4.5\text{ 浬}
\]
答：船与灯塔的距离是 \(4.5\text{ 浬}\).
\end{solution}

\begin{problem}
要在墙上开一矩形的大玻璃窗，周长限定为 \(6\text{ 米}\).

\begin{enumerate}
\item 求以矩形的一条边长 \(x\) 表示窗户的面积 \(y\) 的函数式.
\item 求这函数图象的顶点的坐标，对称轴的方程.
\item 画出这函数的图象，并求出 \(x\) 的允许值的范围.
\end{enumerate}
\end{problem}

\begin{solution}
\begin{enumerate}
\item
设矩形的另一条边长为 \(z\text{ 米}\). 由周长为 \(6\text{ 米}\)，有 \(2x+2z=6\)，所以 \(z=3-x\). 因而面积为
\[
y=xz=x(3-x)=-x^{2}+3x
\]

\item
配方得
\[
y=-x^{2}+3x=-\left(x^{2}-3x+\frac94\right)+\frac94=-\left(x-\frac32\right)^{2}+\frac94
\]
因此函数图象的顶点为 \(\left(\frac32,\frac94\right)\)，对称轴方程为 \(x=\frac32\).

\item
由两条边长均为正数，得 \(x>0\) 且 \(z=3-x>0\)，所以 \(x\) 的允许值范围为 \(0<x<3\). 图象是抛物线 \(y=-x^{2}+3x\) 在该区间内的部分；它与横轴的延长交于 \((0,0)\)、\((3,0)\)，顶点为 \(\left(\frac32,\frac94\right)\)，开口向下，见图.
\end{enumerate}

\solutionfigure{\bitmapfigure[max width=0.36\linewidth]{img/1960/national/q4_solution_fig1.png}}
\end{solution}

\section{选作题}

\begin{problem}
已知方程 \(2x^{2}-x\cdot4\sin\theta+3\cos\theta=0\) 的两个根相等，且 \(\theta\) 为锐角，求 \(\theta\) 和这个方程的两个根.
\end{problem}

\begin{solution}
二次项系数为 \(2\neq0\)，所以方程有两个相等的根当且仅当判别式为零. 因而
\[
\Delta=(-4\sin\theta)^2-4\cdot2\cdot3\cos\theta=0
\]
即 \(2\sin^2\theta-3\cos\theta=0\). 用 \(\sin^2\theta=1-\cos^2\theta\) 化为
\[
2\cos^2\theta+3\cos\theta-2=0
\Longleftrightarrow(2\cos\theta-1)(\cos\theta+2)=0
\]
因 \(-1<\cos\theta<1\)，\(\cos\theta+2=0\) 不可能；故 \(\cos\theta=\frac12\). 又 \(\theta\) 为锐角，所以 \(\theta=60^{\circ}\).

方程的重根为
\[
x_0=-\frac{-4\sin\theta}{2\cdot2}=\sin\theta
\]
所以 \(x_0=\sin60^{\circ}=\frac{\sqrt3}{2}\). 答：\(\theta=60^{\circ}\)，方程的两个根都是 \(\frac{\sqrt3}{2}\).
\end{solution}

\begin{problem}
\(a\) 是什么实数的时候，下列方程组的解是正数：
\(
\begin{cases}
2x+ay=4,\\
x+4y=8
\end{cases}
\)
\end{problem}

\begin{solution}
先考察 \(a=8\). 此时将第二式乘以 \(2\) 得 \(2x+8y=16\)，而第一式为 \(2x+8y=4\)，两式矛盾，所以无解.

当 \(a\neq8\) 时，将第二式乘以 \(2\) 减去第一式，得
\((8-a)y=12\)，\(y=\frac{12}{8-a}\).
再由第二式得
\[
x=8-4y=8-\frac{48}{8-a}=\frac{16-8a}{8-a}=\frac{8(2-a)}{8-a}
\]
因为 \(12>0\)，要使 \(y>0\)，必须且只须 \(8-a>0\)，即 \(a<8\). 在此条件下分母 \(8-a\) 为正，要使 \(x>0\)，必须且只须 \(2-a>0\)，即 \(a<2\). 因此必要条件为 \(a<2\).

反之，当 \(a<2\) 时，\(8-a>0\)、\(2-a>0\)，上式给出的 \(x\)，\(y\) 均为正数，且由消元过程满足原方程组. 所以充分性也成立，故解不等式组
\[
\begin{cases}
8-a>0,\\
2-a>0
\end{cases}
\]
得 \(a<2\). 因此，当且仅当 \(a<2\) 时，原方程组的解是正数.
\end{solution}
